\documentclass[a4paper]{article}

\usepackage{amssymb,mathrsfs,amsmath,amscd,amsthm}
\usepackage[mathcal]{euscript}
\usepackage{stmaryrd}
\usepackage[T1]{fontenc}
\usepackage[utf8]{inputenc}
\usepackage[polish]{babel}
\usepackage{enumitem}
\usepackage{graphicx}
\usepackage{titling}
\usepackage{float}
\usepackage{hyperref}

\hypersetup{
    colorlinks=true,
    linkcolor=blue,
    filecolor=magenta,      
    urlcolor=cyan,
    pdftitle={Overleaf Example},
    pdfpagemode=FullScreen,
    }

\setlength\parindent{0pt}

\title{Deklaracja języka}
\author{Marcin Mazur}

\begin{document}

\maketitle
\section{Opis}
Mam zamiar zrobić kopię języka Latte:\\
\url{https://www.mimuw.edu.pl/~ben/Zajecia/Mrj2021/Latte/} \\
Zmiany (w stosunku do wyżej wymienionej gramatyki):
\begin{itemize}
      \item syntax przybliżony do składni pythona
      \item pass-by-value
      \item for loop
      \item read-only variables
      \item break, continue
      \item pozbycie się typu void
      \item podział na deklaracje i statementy
\end{itemize}
Nietypowo działające konstrukcje:
\begin{itemize}
      \item for loop - ostatni element w nawiasach '()' jest dowolnym Stmt,
            który wykona się po każdej iteracji fora. W szczególności oznacza
            to, że można tworzyć bardziej skomplikowane instrukcje w postaci
            bloków (przykłady zawarte w testach).
      \item read-only variables - możliwość przypisania wartości wyłącznie
            w momencie deklaracji, co skutkuje tym, że zmienna 'const',
            która 'nie dostanie' wartości staje się zmienną 'śmieciową'.
            Przez zmienną 'śmieciową' rozumiem, że nie można na nią przypisać żadnej wartości
            ani żadnej wartości z niej odczytać.
      \item incr, decr - są Stmt, i trzeba je zakończyć średnikiem ';'. Co
            skutkuje tym, że niemożliwe są niektóre naturalne konstrukcje
            typu:\\ \texttt{int x = y++ + 10;} (przykłady zawarte w testach).
\end{itemize}

\section{Tabelka}
01 (trzy typy) \\
02 (literały, arytmetyka, porównania) \\
03 (zmienne, przypisanie) \\
04 (print) \\
05 (while, if) \\
06 (funkcje, rekurencja) \\
07 (przez zmienną / przez wartość) \\
08 (zmienne read-only i pętla for) \\
09 (przesłanianie i statyczne wiązanie) \\
10 (obsługa błędów wykonania) \\
11 (funkcje zwracające wartość) \\
16 (1) (break, continue)

Razem: 22pkt

\end{document}

------------------------------------------------------------------------
\begin{verbatim}

\end{verbatim}

------------------------------------------------------------------------\\
IN:
\begin{verbatim}

\end{verbatim}
OUT:
\begin{verbatim}

\end{verbatim}
